\chapter*{Введение}
\addcontentsline{toc}{chapter}{Введение}

Задачи траекторного управления в трехмерном пространстве возникают в беспилотной авиации, мобильной робототехнике и других системах, которые должны двигаться по заданной пространственной кривой с ограничениями по точности, скорости и устойчивости.
На практике этого обычно недостаточно свести к движению по заранее известной траектории.
Важно еще и подобрать такой режим движения, при котором объект проходит маршрут достаточно быстро, но при этом не выходит за допустимые ошибки и не теряет устойчивость.

Для нелинейных динамических систем эта задача особенно чувствительна к качеству модели и способу построения регулятора.
Линейные приближения удобны на этапе первичного анализа, но при заметных изменениях режима движения, кривизны траектории или ограничений объекта их возможностей часто уже не хватает.
Поэтому на первый план выходит сочетание аналитических методов управления и интеллектуальных методов.

Актуальность работы обусловлена растущим применением беспилотных летательных аппаратов в задачах мониторинга, инспекции объектов, картографирования и автономной навигации.
Для подобных систем важна не только способность следовать заданной траектории, но и эффективное использование доступных динамических возможностей объекта.
При этом использование постоянной скорости движения зачастую приводит либо к излишне консервативному режиму управления, либо к снижению качества слежения на участках со сложной геометрией траектории.
В связи с этим представляет практический интерес разработка методов адаптивного выбора скорости движения, учитывающих текущее состояние объекта и особенности маршрута.

В качестве математической основы в работе используется подход к согласованному траекторному управлению по выходу, изложенный в диссертации С.\,А.~Кима~\cite{kim_algorithms}.
Этот подход был выбран как базовый, поскольку он позволяет формализовать задачу слежения за пространственной кривой, ввести ошибки в сопутствующей системе координат и построить регулятор, учитывающий структуру нелинейной модели объекта.

Проблема, на решение которой направлена работа, состоит в ограниченной эффективности выбора параметрической скорости движения в существующем аналитическом контуре управления.
Базовый регулятор обеспечивает устойчивое движение вдоль траектории, однако при фиксированном значении параметрической скорости V* система работает слишком консервативно.
На простых участках траектории объект движется медленнее, чем это допускается динамическими ограничениями, а на участках со сложной геометрией запас по устойчивости, напротив, быстро уменьшается.
В результате возникает необходимость в адаптивном выборе скорости движения, согласованном с текущим состоянием объекта и геометрией траектории.
Из-за этого возникает необходимость в адаптивном выборе скорости движения, согласованном с текущим состоянием объекта и геометрией траектории.

Цель работы состоит в исследовании методов интеллектуального управления траекторным движением для класса нелинейных систем на примере квадрокоптера и в разработке программной реализации, в которой аналитический регулятор траекторного слежения дополняется интеллектуальным модулем выбора параметрической скорости.

Для достижения поставленной цели были сформулированы следующие задачи:
\begin{enumerate}
    \item рассмотреть постановку задачи траекторного управления нелинейным объектом в трехмерном пространстве;
    \item описать математическую модель квадрокоптера и ограничения;
    \item реализовать базовый алгоритм траекторного слежения по выходу;
    \item подготовить процедуру формирования датасета для подбора допустимой скорости движения;
    \item разработать и обучить нейросетевой модуль для прогнозирования значения V^* по текущим признакам состояния;
    \item провести численное моделирование и сравнить базовый и адаптивный режимы управления.
\end{enumerate}

Практическая значимость работы заключается в том, что полученный подход и идею можно использовать для реализации алгоритмов траекторного управления и для сравнения различных способов адаптивного выбора скорости движения.
Кроме того, сама архитектура решения показывает, как можно сочетать аналитический регулятор, отвечающий за устойчивость, и интеллектуальную надстройку, отвечающую за более эффективный режим движения.
